% ── Formal Echo: Feyerabend ───────────────────────────────────
% Appears after Payne's rejection.
% Tone: disruptive, shifting, pattern-breaking.

\begin{formalecho}{Paul Feyerabend}{The Dissolution of Constraint}

There is always a method.

This is what they are taught first, though not always in those
words. It appears instead as a habit: begin here, proceed in this
way, justify each step, arrive at a conclusion that can be
recognized by others who follow the same path.

The method is not merely a tool. It becomes the condition of
intelligibility itself. To speak meaningfully is to speak within
it.

And yet, at every point where something new appears---something not
already contained within the existing structure---the method
hesitates. Not openly. Not as refusal. As adjustment. A refinement
of rules, an extension of permissible operations, a careful
incorporation of what does not quite fit. The system expands, but
always in a way that preserves its identity. What cannot be
incorporated is not rejected outright; it is redescribed until it
becomes admissible.

This is how stability is maintained.

It is also how blindness is produced.

Consider a result that does not conform. Not an error---this is
important---but a configuration that holds under its own
constraints, that does not collapse when examined, that remains
consistent across multiple paths of approach. It is not incoherent.
It is simply incompatible.

The method does not know what to do with such a result.

So it delays. It reframes. It introduces conditions under which
the result might be reconsidered at a later time, when the
surrounding system has shifted enough to absorb it without rupture.
In the meantime, the result exists in a peculiar state. Not false.
Not accepted. Present, but without standing.

This state is often called irrational.

The term is useful precisely because it does not require precise
definition. It marks a boundary without specifying its nature.

\medskip

There are, broadly, three responses available to a system when it
encounters what it cannot immediately incorporate.

It may attempt to eliminate the anomaly by tightening its rules.
This produces clarity at the cost of excluding large regions of
experience. It may attempt to manage the anomaly by assigning it
provisional status, updating as new information arrives. This
allows flexibility, but may never achieve full resolution. Or it
may abandon the demand for a single universal method, allowing
multiple approaches to coexist, even when they cannot be reconciled.

None of these responses is sufficient on its own.

Elimination produces rigidity. Management produces endless
adjustment. Plurality produces fragmentation.

And yet, taken together, they reveal something that none of them
alone can express: that the demand for method---understood as a
fixed set of rules guaranteeing correct outcomes---is itself part
of the problem.

Not because method is useless, but because its scope is limited.

There are configurations that cannot be reached by following a
prescribed sequence of steps, however refined. There are states
that do not emerge from accumulation or correction, but from a
change in the conditions under which the process operates.

These states are not beyond reason.

They are beyond this form of reason.

\medskip

Most systems do not remain long in the space between methods.

They stabilize. They reassert their rules. They convert what was
once disruptive into something manageable, something that can be
taught, repeated, enforced.

And in doing so, they prepare the conditions for the next
disruption.

What is called irrational marks the points at which the system can
no longer proceed as it has. It is not the opposite of reason.

It is the sign that reason, in its current form, has reached its
limit.

\end{formalecho}
